% 11_conclusioni.tex — Capitolo 11: Conclusioni e limiti.
% Documento didattico "Intelligenza artificiale nel progetto supervised_telemedicina".

\chapter{Conclusioni e limiti}\label{cap:conclusioni}

Siamo arrivati alla fine del percorso. In questo capitolo ripercorriamo le
tappe del progetto \progetto{}, raccogliamo ciò che abbiamo imparato,
dichiariamo con onestà i limiti del lavoro e indichiamo le possibili
evoluzioni. L'obiettivo non è celebrare i numeri, ma capire \emph{perché}
questo piccolo progetto è un buon esempio di ingegneria dell'IA.

\section{Il percorso in sintesi}\label{sec:percorso-sintesi}

Il progetto è nato da un'idea semplice: un sistema a regole cliniche
(capitolo~\ref{cap:teacher}) che assegna una classe di rischio a sei
parametri vitali. Per addestrare un modello servivano dati, e i dati sono
stati generati sinteticamente con un generatore a profili e una buffer
zone (capitolo~\ref{cap:dati}). Il modello scelto è un percettrone
multistrato, costruito da zero in numpy (capitolo~\ref{cap:reti}) e
addestrato con discesa del gradiente e backpropagation
(capitolo~\ref{cap:training}). La valutazione ha mostrato un'accuracy di
$0.9851$ sul test congelato, con gli errori concentrati vicino ai confini
(capitolo~\ref{cap:valutazione}).

Il passo successivo è stato capire \emph{cosa} il modello avesse imparato:
la knowledge distillation (capitolo~\ref{cap:distillazione}) ha reso
esplicito che l'MLP imita il teacher rule-based, con confini più morbidi.
Poi è arrivata la parte più importante: la sicurezza
(capitolo~\ref{cap:sicurezza}), con il safety gate a precedenza assoluta,
la soglia di incertezza e il fallback. Infine, la riproducibilità
(capitolo~\ref{cap:riproducibilita}) ha garantito che tutto questo sia
verificabile: stessi seed, stessi hash, stessi numeri, 129 test.

Ogni tappa ha lasciato un insegnamento che le successive hanno riusato: le
regole hanno fornito le etichette, i dati hanno reso possibile il training,
la rete ha imparato a imitare, la valutazione ha misurato il divario, la
distillazione ha spiegato quel divario, la sicurezza lo ha contenuto e la
riproducibilità ha reso tutto verificabile. Nessuna tappa è un'isola: il
progetto è un sistema, non una somma di componenti.

\section{Cosa abbiamo imparato}\label{sec:cosa-imparato}

Tre lezioni attraversano l'intero documento.

La prima: la \textbf{distillazione è un ponte tra trasparenza e
apprendimento}. Le regole sono leggibili ma rigide; la rete è flessibile
ma opaca. Distillare le regole in una rete non è un compromesso: è un
modo per avere entrambe le cose — la conoscenza esplicita del teacher e la
fluidità del modello. Il safety gate resta a garantire che, dove la rete
non arriva, le regole decidano.

La seconda: la \textbf{sicurezza è un vincolo di progetto, non un
ripensamento}. Il safety gate non è stato aggiunto alla fine per
``correggere'' il modello: è il primo pezzo dell'architettura, e il
teacher stesso riusa la sua logica. Il richiamo 1.0 sui casi critici non è
un risultato statistico: è una proprietà del codice. In un sistema clinico,
questa è la differenza tra un esperimento e un prodotto.

La terza: l'\textbf{incertezza è informazione, non un difetto}. La
confidenza softmax non è una probabilità calibrata, ma è un segnale utile:
quando il modello è incerto, il sistema lo dice, e ripiega su regole
esplicite con un motivo. Nascondere l'incertezza sarebbe stato più
elegante e più pericoloso; dichiararla è la scelta corretta.

C'è anche una lezione di metodo, valida per qualunque progetto di machine
learning: \emph{la trasparenza si progetta, non si aggiunge}. Il safety
gate, i metadati di \codice{AgentOutcome}, gli hash di provenienza e i
guard test non sono decorazioni: sono decisioni architetturali prese
all'inizio, che hanno reso possibile spiegare ogni numero di questo
documento. Chi inizia un progetto pensando prima alla trasparenza e poi
alle metriche ottiene, alla fine, entrambe.

\section{Limiti dichiarati}\label{sec:limiti-dichiarati}

Un progetto onesto dichiara i propri limiti. I principali sono cinque.

\begin{itemize}
  \item \textbf{Dataset sintetico}: i dati non sono reali. Il generatore
        produce campioni plausibili da profili clinici, ma un paziente vero
        non è una gaussiana. Tutte le metriche misurano la fedeltà alle
        regole sintetiche, non l'utilità clinica.
  \item \textbf{Test congelato ottimista}: il test congelato è stato
        generato con la buffer zone, che esclude i casi vicini alle soglie.
        Sul test senza buffer l'accuracy scende da $0.9851$ a $0.9324$:
        il numero ufficiale è reale, ma non va estrapolato ai casi limite.
  \item \textbf{Soglie arbitrarie}: le soglie cliniche derivano dal codice
        legacy e non sono state validate clinicamente. Un valore di 141
        mmHg non è qualitativamente diverso da 139, ma le regole lo
        trattano diversamente.
  \item \textbf{Confidenza non calibrata}: la softmax non produce
        probabilità calibrate. La soglia di incertezza a 0.6 è una scelta
        di progetto ragionevole, non una garanzia statistica.
  \item \textbf{Nessuna validazione con medici}: il progetto non è stato
        valutato da personale clinico. Le classi di rischio e i messaggi
        sono ragionevoli, ma nessun medico ha verificato la loro
        appropriatezza.
\end{itemize}

Questi limiti non invalidano il progetto: lo definiscono. È un progetto
didattico, e il suo valore sta proprio nell'essere piccolo, trasparente e
onesto sui propri confini.

\section{Possibili evoluzioni}\label{sec:evoluzioni}

Dai limiti nascono le evoluzioni. Cinque direzioni sono naturali.

\begin{itemize}
  \item \textbf{Dati reali e anonimizzati}: sostituire (o affiancare) il
        dataset sintetico con dati reali de-identificati, mantenendo la
        pipeline di riproducibilità già costruita.
  \item \textbf{Calibrazione delle probabilità}: trasformare la confidenza
        softmax in una probabilità calibrata (per esempio con temperature
        scaling o Platt scaling), rendendo la soglia di incertezza più
        interpretabile.
  \item \textbf{Explainability}: aggiungere strumenti di spiegazione
        locale come SHAP o LIME per capire \emph{perché} il modello ha
        scelto una classe, accanto al motivo già esplicito del fallback.
  \item \textbf{Monitoraggio in produzione}: tracciare nel tempo la
        distribuzione delle confidenze e dei fallback per rilevare
        \emph{drift} dei dati o degrado del modello.
  \item \textbf{Estensione del compito}: più classi di rischio, o un
        modello di regressione per un indice di gravità continuo, con lo
        stesso schema gate + modello + fallback.
\end{itemize}

\section{Come approfondire}\label{sec:approfondire}

Per chi vuole andare oltre, la bibliografia di questo documento è il punto
di partenza giusto. Il testo di riferimento per le reti profonde è
\emph{Deep Learning} di Goodfellow, Bengio e Courville
\cite{goodfellow2016deep}; per il machine learning in generale, il
\emph{Pattern Recognition and Machine Learning} di Bishop
\cite{bishop2006pattern} e \emph{The Elements of Statistical Learning} di
Hastie, Tibshirani e Friedman \cite{hastie2009elements}. La knowledge
distillation è introdotta nell'articolo originale di Hinton, Vinyals e
Dean \cite{hinton2015distilling}. Tutti i dettagli implementativi del
progetto sono nel repository e nella documentazione tecnica
\cite{progetto2026}.

Il consiglio pratico, per chi vuole verificare da sé: ripercorrere i
capitoli in ordine, eseguendo i comandi del progetto a ogni tappa. Prima
le regole e il generatore, poi il training e la valutazione, poi l'agente
con il safety gate, infine la dashboard. Ogni numero di questo documento
dovrebbe ricomparire identico: è il modo migliore per trasformare la
lettura in comprensione.

\medskip

Il progetto \progetto{} è piccolo: 323 parametri, un MLP scritto a mano in
numpy, un dataset sintetico. Eppure contiene tutto ciò che serve per capire
l'IA applicata: dati, modello, training, valutazione, distillazione,
sicurezza, incertezza e riproducibilità. La lezione finale è questa: un
sistema piccolo, trasparente e sicuro vale più di un modello grande e
opaco. La potenza non sta nel numero di parametri, ma nella capacità di
spiegare ogni decisione, di contenere ogni errore e di riprodurre ogni
risultato. Questo è il senso del percorso che abbiamo fatto insieme.